El presente capítulo define las políticas, responsabilidades y procedimientos establecidos por SKYFILMS para la gestión, registro, control y custodia del Libro de Mantenimiento de cada Sistema de Aeronave No Tripulada (UAS) perteneciente a la organización.

El Libro de Mantenimiento constituye el documento técnico oficial que registra de forma cronológica todas las inspecciones, reparaciones, reemplazos de componentes, discrepancias y verificaciones efectuadas a cada aeronave, garantizando la trazabilidad y evidencia documental de la aeronavegabilidad continuada.

\section{Alcance}
Este procedimiento aplica a todos los UAS operados por SKYFILMS, tanto de propiedad directa como aquellos bajo contrato de operación autorizado, e incluye:

\begin{itemize}
    \item Registro de mantenimiento preventivo y correctivo.

    \item Inspecciones periódicas y verificaciones funcionales.

    \item Sustitución de componentes, accesorios o baterías.

    \item Reporte y cierre de discrepancias técnicas.

    \item Certificación de liberación al servicio.
\end{itemize}

\section{Estructura del Libro de Mantenimiento}
Cada \textbf{Libro de Mantenimiento UAS} debe contener, como mínimo, las siguientes secciones:

\begin{enumerate}
    \item \textbf{Portada de identificación:}
        \begin{itemize}
            \item Nombre del explotador (SKYFILMS).

            \item Modelo y número de serie del UAS.

            \item Fecha de inicio del registro y número de control interno.
        \end{itemize}

    \item \textbf{Registro de mantenimiento menor:}
        \begin{itemize}
            \item Actividades realizadas conforme al Programa de Mantenimiento.

            \item Fecha, horas de vuelo acumuladas, y firma del técnico responsable.
        \end{itemize}

    \item \textbf{Registro de mantenimiento mayor:}    
        \begin{itemize}
            \item Descripción detallada del o los componentes a los cuales se les hará el mantenimiento y descripción del estado.
        \end{itemize}

    \item \textbf{Registro de mantenimiento correctivo:}
        \begin{itemize}
            \item Descripción detallada de la falla, acción correctiva, piezas reemplazadas y resultados de prueba.
        \end{itemize}

    \item \textbf{Liberación para servicio:}
        \begin{itemize}
            \item Firma del \textbf{Técnico UAS} para el retorno a operación.
        \end{itemize}


    \item \textbf{Observaciones y discrepancias reportadas:}
        \begin{itemize}
            \item Notas o eventos técnicos informados por los pilotos o el personal operativo.
        \end{itemize}
\end{enumerate}

Cada formato está numerado correlativamente y se gestiona bajo el control documental establecido por el Técnico de UAS.

\section{Procedimiento de Diligenciamiento}

\begin{itemize}
    \item El \textbf{Técnico UAS} deberá registrar inmediatamente cada acción de mantenimiento en el Libro correspondiente, evitando tachaduras o enmendaduras.

    \item Los registros deberán ser \textbf{claros, legibles y completos}, describiendo la tarea realizada, la referencia del procedimiento técnico o manual del fabricante, y las horas totales de operación.

    \item Las correcciones sólo podrán efectuarse mediante \textbf{anotación complementaria, firmada y fechada}, sin alterar el registro original.

    \item La \textbf{liberación al servicio} de la aeronave sólo podrá registrarse una vez verificada la correcta ejecución del mantenimiento por el Tecnico UAS.

    \item No se permitirá el uso de formularios o registros técnicos no aprobados por SKYFILMS o no controlados por número de versión.
\end{itemize}

\section{Control y Custodia}
\begin{itemize}
    \item El \textbf{Técnico UAS} es el responsable de la custodia física y digital del Libro de Mantenimiento de cada aeronave.

    \item Los Libros activos se conservarán en el área técnica y se mantendrán \textbf{bajo control de acceso restringido}.

    \item Al completar un libro o reemplazarlo por uno nuevo, el anterior se archivará en la \textbf{Oficina Técnica de Mantenimiento}, conservándose por un periodo mínimo de cinco (5) años.

    \item Los libros se mantendrán disponibles para \textbf{revisión o auditoría} por parte de la autoridad aeronáutica en cualquier momento.
\end{itemize}

\section{Registros Electrónicos}
SKYFILMS empleara formatos \textbf{electrónicos} para la gestión del \textbf{Libro de Mantenimiento}, cumpliendo con los siguientes requisitos:

\begin{itemize}
    \item Garantiza la integridad, autenticidad y disponibilidad de los registros.

    \item Cuenta con mecanismos de seguridad y respaldo automático de la información.

    \item Permite la impresión y verificación física de los registros cuando lo solicite la Aerocivil.

    \item Esta aprobados por el Tecnico UAS antes de su implementación.
\end{itemize}

El sistema digital incluye controles de usuario y trazabilidad que permiten identificar al responsable de cada registro o modificación.

\section{Responsabilidades}
\begin{itemize}
    \item El \textbf{Técnico de UAS} es:
        \begin{itemize}
            \item Responsable de la supervisión general de la gestión del Libro de Mantenimiento y de la integridad de los registros.

            \item Responsable controlar la emisión, numeración, custodia y archivo de los libros y formatos.

            \item Responsable del registro oportuno, veraz y completo de cada actividad realizada.
        \end{itemize}

    \item El \textbf{Gerente General} debe garantizar la disponibilidad de recursos y sistemas que aseguren la correcta conservación y acceso a los registros técnicos.
\end{itemize}

\section{Disposiciones Finales}

\begin{itemize}
    \item Ninguna aeronave podrá ser liberada para operación si no cuenta con su \textbf{Libro de Mantenimiento actualizado y debidamente diligenciado}.

    \item La pérdida o deterioro de un libro deberá ser reportada de inmediato al \textbf{Gerente General} y documentada mediante un informe técnico.

    \item El incumplimiento de las disposiciones establecidas en este capítulo será considerado una \textbf{no conformidad técnica grave, sujeta a las acciones correctivas correspondientes}.
\end{itemize}

El Libro de Mantenimiento constituye un elemento esencial para demostrar la aeronavegabilidad continuada y el cumplimiento técnico de la flota UAS de SKYFILMS ante la Unidad Administrativa Especial de Aeronáutica Civil (Aerocivil).